\documentclass[11pt,a4paper]{article}

% ---- Packages ----
\usepackage[utf8]{inputenc}
\usepackage[T1]{fontenc}
\usepackage{amsmath,amssymb,amsthm}
\usepackage{mathtools}
\usepackage[margin=1in,headheight=14pt]{geometry}
\usepackage{hyperref}
\usepackage{enumitem}
\usepackage{xcolor}
\usepackage{tcolorbox}
\usepackage{fancyhdr}
\usepackage{booktabs}

% ---- Hyperref ----
\hypersetup{
    colorlinks=true,
    linkcolor=red,
    citecolor=blue,
    urlcolor=blue
}

% ---- Header ----
\pagestyle{fancy}
\fancyhf{}
\fancyhead[L]{\textit{Probability Foundations for Reinforcement Learning}}
\fancyhead[R]{\thepage}
\renewcommand{\headrulewidth}{0.4pt}

% ---- Theorem Environments ----
\theoremstyle{definition}
\newtheorem{definition}{Definition}[section]
\newtheorem{assumption}{Assumption}[section]
\newtheorem{example}{Example}[section]

\theoremstyle{plain}
\newtheorem{theorem}{Theorem}[section]
\newtheorem{lemma}[theorem]{Lemma}
\newtheorem{proposition}[theorem]{Proposition}
\newtheorem{corollary}[theorem]{Corollary}

\theoremstyle{remark}
\newtheorem{remark}{Remark}[section]

% ---- Colored Boxes ----
\tcbuselibrary{theorems,skins,breakable}

\newtcolorbox{keyresult}[1][]{
    colback=green!5!white,
    colframe=green!50!black,
    fonttitle=\bfseries,
    title=#1,
    breakable
}

\newtcolorbox{rlconnection}[1][]{
    colback=blue!5!white,
    colframe=blue!50!black,
    fonttitle=\bfseries,
    title=#1,
    breakable
}

\newtcolorbox{tdconnection}[1][]{
    colback=blue!5!white,
    colframe=blue!50!black,
    fonttitle=\bfseries,
    title=#1,
    breakable
}

% ---- Operators ----
\DeclareMathOperator{\Var}{Var}
\DeclareMathOperator{\Cov}{Cov}
\newcommand{\E}{\mathbb{E}}
\newcommand{\R}{\mathbb{R}}
\newcommand{\N}{\mathbb{N}}
\newcommand{\Prob}{\mathbb{P}}
\newcommand{\indicator}{\mathbf{1}}
\newcommand{\cas}{\xrightarrow{\text{a.s.}}}
\newcommand{\cprob}{\xrightarrow{\Prob}}
\newcommand{\clp}[1]{\xrightarrow{L^{#1}}}
\newcommand{\cdist}{\xrightarrow{d}}

% ---- Title ----
\title{Lesson 7: Variance, Covariance, and Independence}
\author{AI/ML Mathematics Series}
\date{February 2026}

\begin{document}
\maketitle
\tableofcontents
\newpage

\setcounter{section}{5}

\section{Variance and Conditional Variance}
\label{sec:variance}

\subsection{Definitions}
\label{subsec:var_def}

\begin{definition}[Variance]
\label{def:variance}
The \emph{variance} of a random variable $X$ with $\E[X^2] < \infty$ is
\begin{equation}
    \Var(X) = \E\!\left[(X - \E[X])^2\right].
    \label{eq:var_def}
\end{equation}
\end{definition}

\begin{theorem}[Computational Formula for Variance]
\label{thm:var_formula}
\begin{equation}
    \Var(X) = \E[X^2] - (\E[X])^2.
    \label{eq:var_formula}
\end{equation}
\end{theorem}

\begin{proof}
Let $\mu = \E[X]$. Expanding the square:
\begin{align}
    \Var(X) &= \E[(X - \mu)^2] \nonumber\\
    &= \E[X^2 - 2\mu X + \mu^2] \nonumber\\
    &= \E[X^2] - 2\mu\,\E[X] + \mu^2 \tag{linearity} \\
    &= \E[X^2] - 2\mu^2 + \mu^2 \nonumber\\
    &= \E[X^2] - \mu^2 = \E[X^2] - (\E[X])^2. \label{eq:var_formula_proof}
\end{align}
\end{proof}

\begin{definition}[Covariance]
\label{def:covariance}
The \emph{covariance} of random variables $X$ and $Y$ with $\E[X^2], \E[Y^2] < \infty$ is
\begin{equation}
    \Cov(X, Y) = \E\!\left[(X - \E[X])(Y - \E[Y])\right] = \E[XY] - \E[X]\,\E[Y].
    \label{eq:covariance_def}
\end{equation}
\end{definition}

\begin{proof}[Proof of the computational formula]
Expand the product $(X - \E[X])(Y - \E[Y]) = XY - X\,\E[Y] - \E[X]\,Y + \E[X]\,\E[Y]$.
Taking expectations and using linearity:
$\E[XY] - \E[X]\,\E[Y] - \E[X]\,\E[Y] + \E[X]\,\E[Y] = \E[XY] - \E[X]\,\E[Y]$.
\end{proof}

\begin{remark}
$\Var(X) = \Cov(X, X)$, and $\Cov(X, Y) = 0$ when $X$ and $Y$ are independent (by Theorem~\ref{thm:product_exp} below, $\E[XY] = \E[X]\,\E[Y]$).
\end{remark}

\begin{definition}[Conditional Variance]
\label{def:cond_var}
The \emph{conditional variance} of $X$ given $Y$ is the random variable
\begin{equation}
    \Var(X \mid Y) = \E\!\left[(X - \E[X \mid Y])^2 \;\middle|\; Y\right] = \E[X^2 \mid Y] - (\E[X \mid Y])^2.
    \label{eq:cond_var}
\end{equation}
\end{definition}

\begin{proof}[Proof of the computational formula~\eqref{eq:cond_var}]
Fix $Y = y$ and let $\mu_y = \E[X \mid Y = y]$. Applying the same expansion as in Theorem~\ref{thm:var_formula} to the conditional expectation:
\begin{align}
    \Var(X \mid Y = y) &= \E[(X - \mu_y)^2 \mid Y = y] \nonumber\\
    &= \E[X^2 \mid Y = y] - 2\mu_y\,\E[X \mid Y = y] + \mu_y^2 \nonumber\\
    &= \E[X^2 \mid Y = y] - \mu_y^2. \label{eq:cond_var_proof}
\end{align}
\end{proof}

\subsection{Law of Total Variance}
\label{subsec:total_var}

\begin{keyresult}[Law of Total Variance]
\begin{theorem}[Law of Total Variance]
\label{thm:total_variance}
For random variables $X$ and $Y$ with $\E[X^2] < \infty$:
\begin{equation}
    \Var(X) = \E\!\left[\Var(X \mid Y)\right] + \Var\!\left(\E[X \mid Y]\right).
    \label{eq:total_variance}
\end{equation}
\end{theorem}
\end{keyresult}

\begin{proof}
We compute each term on the right-hand side.

\textbf{First term.} Using~\eqref{eq:cond_var}:
\begin{align}
    \E[\Var(X \mid Y)]
    &= \E\!\left[\E[X^2 \mid Y] - (\E[X \mid Y])^2\right] \nonumber\\
    &= \E[\E[X^2 \mid Y]] - \E[(\E[X \mid Y])^2] \nonumber\\
    &= \E[X^2] - \E[(\E[X \mid Y])^2], \label{eq:total_var_first}
\end{align}
where the last equality uses the tower property $\E[\E[X^2 \mid Y]] = \E[X^2]$.

\textbf{Second term.} Let $\mu = \E[X]$. By the computational formula for variance:
\begin{align}
    \Var(\E[X \mid Y])
    &= \E[(\E[X \mid Y])^2] - (\E[\E[X \mid Y]])^2 \nonumber\\
    &= \E[(\E[X \mid Y])^2] - (\E[X])^2 \nonumber\\
    &= \E[(\E[X \mid Y])^2] - \mu^2. \label{eq:total_var_second}
\end{align}

\textbf{Sum.} Adding~\eqref{eq:total_var_first} and~\eqref{eq:total_var_second}:
\begin{align}
    \E[\Var(X \mid Y)] + \Var(\E[X \mid Y])
    &= \bigl(\E[X^2] - \E[(\E[X \mid Y])^2]\bigr) + \bigl(\E[(\E[X \mid Y])^2] - \mu^2\bigr) \nonumber\\
    &= \E[X^2] - \mu^2 = \Var(X). \label{eq:total_var_sum}
\end{align}
\end{proof}

\begin{tdconnection}[Connection to TD Learning]
This is the central result for the bias-variance analysis. The article applies it with $X = G_t$ and $Y = (A_t, S_{t+1})$:
\[
    \Var_\pi(G_t \mid S_t) = \underbrace{\Var_\pi(R_{t+1} + \gamma V^\pi(S_{t+1}) \mid S_t)}_{\text{one-step variance}} + \gamma^2 \underbrace{\E_\pi[\Var_\pi(G_{t+1} \mid S_{t+1}) \mid S_t]}_{\text{future variance} \;\geq\; 0}.
\]
The TD(0) target $R_{t+1} + \gamma V(S_{t+1})$ captures only the first term, eliminating the non-negative future variance entirely. This is why TD has lower variance than Monte Carlo.
\end{tdconnection}

%% ============================================================
%% SECTION 7: INDEPENDENCE
%% ============================================================

\section{Independence}
\label{sec:independence}

\subsection{Independent Events}
\label{subsec:indep_events}

\begin{definition}[Independent Events]
\label{def:indep_events}
Two events $A, B \in \mathcal{F}$ are \emph{independent} (written $A \perp B$) if:
\begin{equation}
    \Prob(A \cap B) = \Prob(A)\, \Prob(B).
    \label{eq:indep_events}
\end{equation}
Events $A_1, \ldots, A_n$ are \emph{mutually independent} if for every subset $I \subseteq \{1, \ldots, n\}$:
\begin{equation}
    \Prob\!\left(\bigcap_{i \in I} A_i\right) = \prod_{i \in I} \Prob(A_i).
    \label{eq:mutual_indep}
\end{equation}
\end{definition}

\begin{remark}
Mutual independence is strictly stronger than pairwise independence. Pairwise independence requires~\eqref{eq:indep_events} for every pair $(A_i, A_j)$ with $i \neq j$, but this does not imply~\eqref{eq:mutual_indep} for $|I| \geq 3$.
\end{remark}

\begin{example}[Pairwise but Not Mutual Independence]
\label{ex:pairwise_not_mutual}
Let $\Omega = \{1, 2, 3, 4\}$ with uniform probability $1/4$ each. Define $A = \{1,2\}$, $B = \{1,3\}$, $C = \{1,4\}$. Then $\Prob(A) = \Prob(B) = \Prob(C) = 1/2$ and $\Prob(A \cap B) = \Prob(A \cap C) = \Prob(B \cap C) = 1/4$, so all pairs are independent. But $\Prob(A \cap B \cap C) = \Prob(\{1\}) = 1/4 \neq 1/8 = \Prob(A)\Prob(B)\Prob(C)$.
\end{example}

\subsection{Independent Random Variables}
\label{subsec:indep_rv}

\begin{definition}[Independent Random Variables]
\label{def:indep_rv}
Discrete random variables $X_1, \ldots, X_n$ are \emph{mutually independent} if for all $(x_1, \ldots, x_n)$:
\begin{equation}
    \Prob(X_1 = x_1, \ldots, X_n = x_n) = \prod_{i=1}^n \Prob(X_i = x_i).
    \label{eq:indep_rv}
\end{equation}
Equivalently, the joint PMF factors as a product of marginals.
\end{definition}

\begin{keyresult}[Product of Expectations]
\begin{theorem}[Expectation of Product of Independent Random Variables]
\label{thm:product_exp}
If $X_1, \ldots, X_n$ are mutually independent with $\E[|X_i|] < \infty$, then:
\begin{equation}
    \E\!\left[\prod_{i=1}^n X_i\right] = \prod_{i=1}^n \E[X_i].
    \label{eq:product_exp}
\end{equation}
\end{theorem}
\end{keyresult}

\begin{proof}
We prove the case $n = 2$; the general case follows by induction. Let $X, Y$ be independent discrete random variables. Then:
\begin{align}
    \E[XY] &= \sum_x \sum_y xy\, \Prob(X = x, Y = y) \nonumber\\
    &= \sum_x \sum_y xy\, \Prob(X = x)\, \Prob(Y = y) \tag{independence}\\
    &= \left(\sum_x x\, \Prob(X = x)\right)\left(\sum_y y\, \Prob(Y = y)\right) = \E[X]\, \E[Y]. \label{eq:product_exp_proof}
\end{align}
The factoring of the double sum into a product of single sums uses $xy = x \cdot y$ and the distributive law, which is valid when both sums converge absolutely.
\end{proof}

\begin{corollary}[Variance of Sum of Independent Random Variables]
\label{cor:var_sum}
If $X_1, \ldots, X_n$ are mutually independent with $\E[X_i^2] < \infty$, then:
\begin{equation}
    \Var\!\left(\sum_{i=1}^n X_i\right) = \sum_{i=1}^n \Var(X_i).
    \label{eq:var_sum}
\end{equation}
\end{corollary}

\begin{proof}
Let $S = \sum_i X_i$ and $\mu_i = \E[X_i]$. Then:
\begin{align}
    \Var(S) &= \E\!\left[\left(\sum_i (X_i - \mu_i)\right)^2\right] = \sum_i \E[(X_i - \mu_i)^2] + 2 \sum_{i < j} \E[(X_i - \mu_i)(X_j - \mu_j)].
\end{align}
For independent $X_i, X_j$, the centered variables $X_i - \mu_i$ and $X_j - \mu_j$ are also independent, so by Theorem~\ref{thm:product_exp}:
\begin{equation}
    \E[(X_i - \mu_i)(X_j - \mu_j)] = \E[X_i - \mu_i]\, \E[X_j - \mu_j] = 0 \cdot 0 = 0.
\end{equation}
Thus $\Var(S) = \sum_i \Var(X_i)$.
\end{proof}

\subsection{Conditional Independence}
\label{subsec:cond_indep}

\begin{definition}[Conditional Independence]
\label{def:cond_indep}
Events $A$ and $B$ are \emph{conditionally independent given $C$} (with $\Prob(C) > 0$) if:
\begin{equation}
    \Prob(A \cap B \mid C) = \Prob(A \mid C)\, \Prob(B \mid C).
    \label{eq:cond_indep}
\end{equation}
Random variables $X$ and $Y$ are \emph{conditionally independent given $Z$} if for all $x, y, z$ with $\Prob(Z = z) > 0$:
\begin{equation}
    \Prob(X = x, Y = y \mid Z = z) = \Prob(X = x \mid Z = z)\, \Prob(Y = y \mid Z = z).
    \label{eq:cond_indep_rv}
\end{equation}
\end{definition}

\begin{remark}
Independence and conditional independence are \emph{distinct} properties. Independent random variables may become conditionally dependent (and vice versa). In particular, conditional independence given $Z$ does \emph{not} imply unconditional independence.
\end{remark}

\begin{rlconnection}[Connection to RL]
The \textbf{Markov property} is a conditional independence statement: given $S_t$, the future $(S_{t+1}, R_{t+1}, S_{t+2}, \ldots)$ is conditionally independent of the past $(S_0, A_0, \ldots, S_{t-1}, A_{t-1})$. This is the foundational structural assumption of MDPs and is used in every Bellman equation derivation and every TD convergence proof.
\end{rlconnection}

%% ============================================================
%% SECTION 8: FILTRATIONS
%% ============================================================


\end{document}